\chapter{Introduction}\label{ch:introduction}

BETA AID is a customer-service copilot for LATAM Bank, built for the Factored AI \& Data Hackathon 2026, together with
the data platform beneath it: a lakehouse of bronze, silver and gold layers, orchestrated pipelines, data contracts,
an audit layer, and the readiness gates that decide which machine-learning products the data can support. Before any
of it was built, a long series of exploratory notebooks asked what the bank's data contain, whether they can be
trusted, at which grain they carry information, and which targets can be learned from them. The Grain Atlas, the
platform's control plane, publishes their answers. This report examines the methods behind those answers: what each
one assumes, whether the assumption holds in these data, how much uncertainty its result carries, and what follows
for the pipelines and the gates.

\section{The data}\label{sec:intro-data}

The bank operates in Mexico, Colombia and Argentina. Its dataset was delivered as a set of extracts with a data
dictionary and a summary of target volumes: customers, their products, transactions, digital events, contact-centre
interactions and their transcripts, complaints, satisfaction surveys, marketing campaigns and their sends, branches,
service agents, and daily exchange rates. The main copy holds $\est{ch03.tables-main}{value}$ tables, among them
$\est{ch03.rows-main-customers}{value}$ customers and $\est{ch03.rows-main-transactions}{value}$ transactions over
$\est{ch03.main-tx-months}{value}$ months, and ends on \proto{dataset.end_date}.

A second copy arrived with it, which no document mentions: the backup. It holds $\est{ch03.tables-backup}{value}$
tables, and its transactions end on \proto{dataset.backup_transactions_end_date}, after
$\est{ch03.backup-tx-months}{value}$ months. The exploratory series set out hypotheses about how the copies relate, and
Part~II tests them.
\begin{description}[style=nextline]
  \item[\ident{H1}, a shift in time] The backup holds the same records on other dates (\cref{ch:time-and-clocks}).
  \item[\ident{H2}, new keys] It holds the same entities under new surrogate keys (\cref{ch:inventory-keys}).
  \item[\ident{H3}, an independent replicate] It is another draw from the same generating process
  (\cref{ch:statistical-equivalence}).
  \item[\ident{H4}, an incomplete copy] It lacks tables or periods of the main copy (\cref{ch:inventory-keys}).
  \item[\ident{H5}, silent mutation] Records present in both copies were changed without a trace
  (\cref{ch:record-linkage}).
\end{description}
The answer matters beyond the backup. If the copies are independent draws of one process, the backup is a replicate
against which a finding on the main copy can be checked, and the report uses it that way.

\section{The readiness question}\label{sec:intro-question}

A machine-learning product is worth building only when the data carry a signal that a model can learn out of time,
that is larger than what a transparent rule already captures, and that holds across markets. The project turns that
requirement into readiness gates. For each candidate model, a gate compares the model with a transparent benchmark on
later data and returns green (build it), amber (use the signal as a rule or a KPI) or red (the data do not support
it), together with the root cause and the data requirement that would turn it green (\cref{sec:em-gate}). The
published scorecard of the $\est{ch19.gates}{value}$ gates of the hour grain has $\est{ch19.pub-green}{value}$ green,
$\est{ch19.pub-amber}{value}$ amber and $\est{ch19.pub-red}{value}$ red gates.

A red gate is not a dead end. It is an audit: it says whether a target is unlearnable because of the natural structure
of the data, a defect of the data, the wrong capture or grain, or missing complementary data, and each cause has an
owner and a fix in the data-engineering cycle. Read this way, the gates are the measuring instrument of a process that
turns unlearnable data into learnable data, and this report is the calibration of that instrument.

\section{Why a report}\label{sec:intro-why}

The notebooks are exploratory. They chose features, windows, benchmarks and tests after looking at the data; they
computed many intervals as if every row were independent, although a customer's transactions, an agent's contacts and
the days of a calendar are not; and they ran many tests without controlling for multiplicity. Such analyses can make
chance findings look real \autocite{simmons2011}. A gate is a decision with consequences, so its evidence should be the
strictest the data allow. The report therefore re-examines every method with the tools that \cref{ch:evidence-model}
develops: intervals computed over the unit that is actually independent, false-discovery control over families of tests
fixed in advance, and a registered re-estimation of every gate under a specification frozen before it ran
\autocite{nosek2018}.

The report also enforces a discipline on itself. No number derived from data is typed into the text: each is a field of
a row of a claims ledger, generated from the artefact it was computed from, whose digest the row records. A claim's
status says what it may be used for: exploratory, registered or positive control. The test suite, which runs on every
commit, fails on a typed number, a
missing claim, a stale digest or an unmapped notebook heading. The chapters of Parts~II to~VI rest on
$\est{ch01.ledger-claims}{value}$ claims (\cref{app:claims-ledger}), and the $\est{ch01.ledger-headings}{value}$
headings of the $\est{ch01.ledger-notebooks}{value}$ notebook sources are mapped to the chapters that treat them
(\cref{app:traceability}).

\section{Contributions}\label{sec:intro-contributions}

\begin{enumerate}
  \item \textbf{An evidence model for nested, serially dependent data} (\cref{ch:evidence-model}): cluster, wild cluster
  and stationary block bootstraps chosen by the number and kind of independent units, false-discovery control over
  declared families, a claim status that limits what each number may decide, and a registered re-estimation whose
  limits are stated as plainly as its guarantees.
  \item \textbf{An account of the data as delivered} (Part~II): what each copy contains and how the copies relate, in
  which clock every process drew its days, which records were silently changed, whether the copies are draws of the
  same process, and what the source changed without saying so, with the controls that would and would not notice.
  \item \textbf{Calibrated monitoring} (Part~III): for every detector, rule and chart, the size of effect it would have
  caught and the rate at which it alarms when nothing happens, so that a null result is informative.
  \item \textbf{The pipeline as a statistical object} (Part~IV): the assumptions embedded in conformance, features,
  marts, privacy and audit, checked against the built lakehouse, with the defects found and their repairs.
  \item \textbf{The grain at which information lives} (Part~V): the rules of aggregation, and the customer, operational
  and campaign signals that survive them.
  \item \textbf{Learnability and readiness} (Part~VI): the learnability map of $\est{ch18.targets}{value}$ targets in
  $\est{ch18.scopes}{value}$ scopes, a calibrated model with a pair of parameters for the one learnable target, and
  the registered re-estimation of every gate on both copies.
\end{enumerate}

\section{How to read this report}\label{sec:intro-reading}

Part~I sets out the evidence model, and every later chapter relies on it. Parts~II to~V follow the pipeline from the
raw extracts to the grains of the gold layer, and Part~VI turns the findings into decisions. Each method section follows
the same order: the question, the estimand, the model and its assumptions stated formally, why this method rather than
another, estimation and inference, diagnostics, results, interpretation, implications for the pipeline and the gates,
and limitations. Each chapter ends with its verdict and its implications, so a reader interested in decisions can read
those sections and \cref{ch:conclusions}.

The conventions before this chapter explain how to read a cited number, its interval and its status. The appendices
give the traceability of every notebook heading (\cref{app:traceability}), the full claims ledger
(\cref{app:claims-ledger}), the country series side by side (\cref{app:per-country}), the software, data digests and
provenance needed to rebuild every number (\cref{app:software}), and the notation (\cref{app:notation}).
